% =============================================================================
% cap07.tex — DIMENSIONAMENTO
% =============================================================================

\chapter{DIMENSIONAMENTO}\label{cap:dimensionamento}

\section{CRITÉRIO DE VERIFICAÇÃO}\label{sec:criterio-verificacao}

Todas as verificações são expressas por uma taxa de utilização, definida como a
razão entre o efeito de cálculo das ações e a resistência de cálculo
correspondente:

\begin{align}
    \eta = \frac{E_d}{R_d} \label{eq:taxa-utilizacao}
\end{align}

\noindent
A verificação da \autoref{eq:taxa-utilizacao} é atendida quando $\eta \leq 1$. Adota-se, além desse limite
normativo, um limite de projeto mais restritivo de $\eta \leq 0{,}75$ para os
modos de ruptura frágeis das ligações, pelas razões expostas no
\autoref{cap:detalhamento}.

O conjunto compreende cinquenta e cinco verificações, das quais
cinquenta e três são executadas e duas declaradas não aplicáveis. As não aplicáveis são o tirante de
base e o tirante de piso de robustez, ambos sem função no modelo de viga
biapoiada adotado na \autoref{ssec:modelo-apoio}: não havendo empuxo a
equilibrar, não há tirante a dimensionar.

Este capítulo apresenta as verificações dos elementos e dos sistemas. As das
ligações e da ancoragem constam do \autoref{cap:detalhamento}. As tabelas
reúnem as taxas de utilização; a verificação governante de cada elemento é
desenvolvida em equações, com a expressão normativa, a substituição dos valores
e o resultado, no \autoref{ap:memoria-calculo}, gerado pelo mesmo programa que
produz os resultados.

\section{VIGA PRINCIPAL}\label{sec:verif-viga}

A \autoref{tab:verif-viga} reúne as verificações de estado-limite último da viga
de duplo-bisel, segundo as cláusulas da EN 1995-1-1 \cite{EN1995-1-1:2004} e com
os esforços de cálculo da \autoref{tab:esforcos-governantes}. O material é GL24h
em classe de serviço 3, com $k_{mod} = 0{,}65$ para ações de duração média.

\tabela{Verificações de estado-limite último da viga principal}{\footnotesize
\begin{tabular}{llcc}
\toprule
    \textbf{Verificação} & \textbf{Cláusula} & \textbf{$\eta$} & \textbf{Situação} \\
\midrule
    Cisalhamento                    & EC5 6.1.7          & 0,3814 & Atende \\
    Tensões na zona de cumeeira     & EC5 6.4.3          & 0,3526 & Atende \\
    Tensões na borda inclinada      & EC5 6.4.2          & 0,3140 & Atende \\
    Compressão no apoio sobre a sela central, L1c & EC5 6.1.5 & 0,4887 & Atende \\
    Compressão no apoio sobre o pilar lateral, L1 & EC5 6.1.5 & 0,4561 & Atende \\
    Flexo-compressão                & EC5 6.3.2 e 6.3.3  & 0,1672 & Atende \\
    Flambagem lateral sob sucção, intradorso comprimido & EC5 6.3.3 & 0,1582 & Atende \\
    Flambagem lateral por torção    & EC5 6.3.3          & 0,1039 & Atende \\
\bottomrule
\end{tabular}}{Elaborado pelo autor (2026)}{verif-viga}{}{}

O cisalhamento é a verificação mais solicitada da viga, com $\eta = 0{,}3814$.
Sua relevância decorre da redução da largura efetiva da seção pelo coeficiente
de fissuração $k_{cr} = 0{,}67$, que retira um terço da largura bruta antes do
cálculo da tensão. Sem essa redução, a taxa cairia para cerca de 0,26.

As tensões na borda inclinada merecem detalhamento, por ilustrarem a assimetria
mencionada na \autoref{sec:vigas-secao-variavel}. Para a inclinação de
10$^{\circ}$, o coeficiente $k_{m,\alpha}$ vale 0,758 quando a borda inclinada
está comprimida e apenas 0,414 quando tracionada. A verificação governante
ocorre a 3,48 m do apoio, sob combinação de sucção, situação em que a borda
superior passa a tracionada. É por essa razão que a verificação não é governada
pela seção de maior momento, e sim por uma seção intermediária sob carregamento
invertido.

A flambagem lateral sob momento positivo apresenta a menor taxa do conjunto, com
$\eta = 0{,}1039$. Esse valor é inteiramente condicionado pelo sistema de
contraventamento: a distância entre contenções do extradorso comprimido, de
2,53 m, corresponde ao espaçamento entre terças, válido apenas porque o
diafragma de cobertura foi verificado como capaz de estabilizar. Ao valor
$0{,}9 \times 2{,}53$ m a nota da tabela 6.1 soma $2h$, por estar a carga
aplicada na borda comprimida, e o comprimento efetivo resulta em 7,85 m.
Caso não fosse, as terças deixariam de ser contenções laterais e o comprimento
efetivo passaria ao vão da própria viga, de 20,25 m. Mesmo nessa situação a taxa
subiria apenas a 0,2026, porque a seção de 0,32 m de largura é robusta à
flambagem lateral, e a verificação continuaria longe de governar.

Nas combinações de sucção o momento inverte, atingindo $-394{,}4$ kN$\cdot$m na
combinação $0{,}9\,G + 1{,}5\,W$, e a borda comprimida passa a ser o intradorso,
que as terças, fixadas no extradorso, não contêm. Essa situação foi verificada
em separado, com hipóteses conservadoras: contenção lateral apenas nos apoios,
sem crédito à contenção da borda tracionada pelas terças, e sem a redução de
$0{,}5h$ que a nota da tabela 6.1 admite quando a carga atua na borda
tracionada. O comprimento efetivo é, então,

\begin{equation}
    \ell_{ef} = 0{,}9\,L = 0{,}9 \times 20{,}25 = 18{,}23~\text{m}.
    \label{eq:lef-succao}
\end{equation}

\noindent
Como a altura da viga varia, a seção crítica foi obtida por varredura ao longo
do vão, nas quinze combinações com trecho de momento negativo. Governa a seção
a 15,80 m do apoio, onde $h = 1{,}786$ m e $M_d = -290{,}4$ kN$\cdot$m, e não a de
maior momento, mais alta. Nessa seção, pelas fórmulas (6.30) a (6.34) da
EN 1995-1-1 \cite{EN1995-1-1:2004},

\begin{equation}
    \sigma_{m,crit} = \frac{0{,}78\,b^2 E_{0,05}}{h\,\ell_{ef}}
    = \frac{0{,}78 \times 0{,}32^2 \times 9.600}{1{,}786 \times 18{,}23}
    = 23{,}56~\text{MPa},
\end{equation}

\noindent
$\lambda_{rel,m} = \sqrt{24/23{,}56} = 1{,}01$ e
$k_{crit} = 1{,}56 - 0{,}75 \times 1{,}01 = 0{,}80$. A tensão de cálculo, de
1,71 MPa, compara-se a $k_{crit} f_{m,d} = 0{,}80 \times 13{,}44 = 10{,}79$ MPa,
com $k_{mod} = 0{,}70$ da ação de curta duração, e resulta $\eta = 0{,}1582$.

\section{PILARES}\label{sec:verif-pilares}

As verificações dos três tipos de pilar, em cada situação de carregamento
pertinente, constam da \autoref{tab:verif-pilares}, segundo as cláusulas da
EN 1995-1-1 \cite{EN1995-1-1:2004}. A seção é de 0,32 $\times$ 0,27 m nos pilares
lateral e central, e de 0,32 $\times$ 0,32 m no de canto.

\tabela{Verificações de estado-limite último dos pilares}{\footnotesize
\begin{tabular}{llcc}
\toprule
    \textbf{Pilar e situação} & \textbf{Cláusula} & \textbf{$\eta$} & \textbf{Situação} \\
\midrule
    De canto, gravidade com contraventamento & EC5 6.2.3 e 6.3.2 & 0,4333 & Atende \\
    Central montante, gravidade com contraventamento & EC5 6.2.3 e 6.3.2 & 0,7460 & Atende \\
    Central, compressão com instabilidade    & EC5 6.3.2         & 0,4151 & Atende \\
    Lateral, vento perpendicular à cumeeira  & EC5 6.3.2         & 0,2587 & Atende \\
    Lateral, vento paralelo à cumeeira       & EC5 6.3.2         & 0,2433 & Atende \\
    Lateral, flexo-compressão                & EC5 6.3.2         & 0,2075 & Atende \\
\bottomrule
\end{tabular}}{Elaborado pelo autor (2026)}{verif-pilares}{}{}

O pilar central recebe esforço normal de cálculo de 199,2 kN, correspondente à
soma das reações de duas vigas, e opera com coeficiente de instabilidade
$k_c = 0{,}445$. Sua taxa de 0,4151 é o dobro da do pilar lateral, com 0,2075,
o que confirma a decisão registrada na \autoref{sec:predim-pilares} de dimensionar
a seção comum pelo pilar central.

O pilar de canto, dotado de seção maior, recebe
esforço normal de 297,0 kN, resultante da reação de gravidade somada ao esforço
descendente de 197,4 kN transmitido pelos dois sistemas de contraventamento que
nele se encontram. Sob sucção, a mesma seção trabalha à flexo-tração com
$\eta = 0{,}224$, e o esforço de arrancamento de 185,3 kN determina a ancoragem
tratada no \autoref{cap:detalhamento}.

O pilar central montante, situado nas duas baías extremas do alinhamento central,
recebe as reações das duas vigas, de 199,2 kN, somadas ao esforço descendente de
129,7 kN do painel de contraventamento, totalizando 328,9 kN. Com a seção comum
de 0,32 $\times$ 0,27 m, resulta $\eta = 0{,}7460$, a maior taxa entre os
pilares.

\section{TERÇAS E COBERTURA}\label{sec:verif-tercas}

A \autoref{tab:verif-tercas} reúne as verificações das terças e dos elementos de
cobertura por elas suportados. As terças são de 0,20 $\times$ 0,58 m em GL24h,
espaçadas de 2,53 m. A verificação da fixação segue a
ABNT NBR 6123 \cite{NBR6123:2023}, com coeficiente de pressão local de
$\pm 2{,}0$; as demais, a EN 1995-1-1 \cite{EN1995-1-1:2004} e a
EN 1991-1-1 \cite{EN1991-1-1:2002}.

\tabela{Verificações de estado-limite último das terças e da cobertura}{\footnotesize
\begin{tabular}{llcc}
\toprule
    \textbf{Verificação} & \textbf{Cláusula} & \textbf{$\eta$} & \textbf{Situação} \\
\midrule
    Fixação da telha sob sucção    & NBR 6123 7.2.6     & 0,6788 & Atende \\
    Terça como biela do diafragma  & EC5 6.3.2          & 0,2634 & Atende \\
    Terça à flexão                 & EC5 6.1.6 e 6.1.7  & 0,0979 & Atende \\
    Terça sob carga concentrada    & EN 1991-1-1 6.3.4.2 & 0,0546 & Atende \\
\bottomrule
\end{tabular}}{Elaborado pelo autor (2026)}{verif-tercas}{}{}

A terça acumula duas funções, e a segunda é a mais exigente. Além de resistir à
flexão devida ao carregamento da cobertura, com taxa de apenas 0,0979, ela
funciona como biela comprimida do diafragma de contraventamento, recebendo
esforço normal de 106,1 kN. Verificada à interação combinada de flexão e
compressão, a taxa sobe para 0,2634, quase o triplo da flexão isolada. A
esbeltez em torno do eixo de menor inércia, com $\lambda_z = 100$, conduz a
$k_{c,z} = 0{,}368$ e explica a diferença.

A fixação da telha é a verificação mais solicitada deste grupo, com
$\eta = 0{,}6788$. A pressão dinâmica de 0,890 kN/m\textsuperscript{2}, associada
ao coeficiente local de sucção de $-2{,}0$, exige fixador a cada 0,20 m ao longo
da terça, isto é, em cada nervura da telha. O espaçamento máximo admissível para
que a taxa não exceda a unidade é de 0,29 m.

%% TODO: dado ausente — a resistência característica do fixador, de 2,5 kN,
%% é valor típico de catálogo e não de ficha técnica. Governa o arrancamento da
%% chapa fina, e não o da madeira, de modo que a confirmação depende da telha
%% efetivamente adotada. Ver checks/elu/fixacao_telha.py.

\section{CONTRAVENTAMENTO E ESTABILIDADE GLOBAL}\label{sec:verif-contraventamento}

Os sistemas que estabilizam a edificação e conduzem as forças horizontais à
fundação são verificados na \autoref{tab:verif-contraventamento}, segundo as
cláusulas da EN 1995-1-1 \cite{EN1995-1-1:2004}, da
EN 1993-1-1 \cite{EN1993-1-1:2005} e da EN 1993-1-8 \cite{EN1993-1-8:2005}. As
diagonais são em aço S235, e as suas ligações de extremidade são detalhadas na
\autoref{sec:detalhe-contraventamento}.

\tabela{Verificações dos sistemas de contraventamento e da estabilidade global}{\footnotesize
\begin{tabular}{llcc}
\toprule
    \textbf{Sistema} & \textbf{Cláusula} & \textbf{$\eta$} & \textbf{Situação} \\
\midrule
    Contraventamento vertical, fachadas & EC5 5.4.4 e EN 1993-1-1 6.2.3 & 0,8352 & Governa o projeto \\
    Contraventamento vertical, alinhamento central & EC5 5.4.4 e EN 1993-1-1 6.2.3 & 0,8352 & Governa o projeto \\
    Contraventamento transversal    & EC5 5.4.4 e EN 1993-1-1 6.2.3 & 0,8242 & Atende \\
    Ligação da diagonal vertical    & EN 1993-1-8 3.13    & 0,8105 & Atende \\
    Ligação da diagonal transversal & EN 1993-1-8 3.13    & 0,8091 & Atende \\
    Contraventamento da cobertura   & EC5 9.2.5           & 0,8009 & Atende \\
    Estabilidade global             & EN 1993-1-1 5.2.1, por analogia & 0,6062 & Atende \\
    Corda do diafragma no beiral    & EC5 6.2.3 e 6.3.2   & 0,0656 & Atende \\
\bottomrule
\end{tabular}}{Elaborado pelo autor (2026)}{verif-contraventamento}{}{}

O contraventamento da cobertura, com $\eta = 0{,}8009$, tem a terceira maior
taxa do projeto. Pelo item 9.2.5.3 da EN 1995-1-1 \cite{EN1995-1-1:2004}, o
sistema resiste à força de estabilização em adição às ações horizontais
externas, de modo que a reação em cada apoio da treliça soma a parcela do vento
paralelo à cumeeira, 104,7 kN, à da estabilização das vigas, 1,4 kN, e vale
106,1 kN; a flecha limite de $\ell/500$ é verificada sob a mesma soma. Como a
força de estabilização decorre da envoltória de todas as combinações, e não da
combinação com vento paralelo à cumeeira, a soma é conservadora. A diagonal de
28 mm recebe esforço de 115,9 kN e
responde por peso de aço da ordem de 1.943 kg, com cruzes nos oito painéis
do vão de cada baía rígida. Como exposto na
\autoref{ssec:o-que-governa}, essa taxa é praticamente irrelevante à seção da
viga, o que a torna o limite efetivo do dimensionamento.

O contraventamento vertical, com $\eta = 0{,}7883$, conduz à fundação a força
longitudinal de 852,5 kN, composta pelos 837,3 kN de vento paralelo à cumeeira e
por 15,2 kN de força de imperfeição. Metade segue aos quatro painéis das
fachadas e metade aos quatro do alinhamento central, na proporção das reações das
treliças da cobertura. Cada painel recebe 106,6 kN, e a diagonal de 33 mm é
solicitada a 167,9 kN. Com extremidades recalcadas, a resistência é a de
escoamento da seção bruta, 201,0 kN, e resulta $\eta = 0{,}8352$, a maior taxa do
projeto. O contraventamento
transversal, menos solicitado, conduz 111,2 kN por meio do diafragma de cobertura
e dos frontões.

As ligações de extremidade das diagonais figuram entre as mais solicitadas do
conjunto, com 0,8105 na diagonal vertical e 0,8091 na transversal. Elas são
dimensionadas por capacidade, isto é, para a resistência da própria barra e não
para o esforço solicitante. Como as barras têm extremidades recalcadas, essa
resistência é a de escoamento, e um eventual excesso de solicitação leva ao
escoamento da barra antes da ruptura do olhal, do pino ou da rosca. O detalhamento consta da \autoref{sec:detalhe-contraventamento}.

A estabilidade global apresenta $\eta = 0{,}6062$, valor que traduz a folga do
fator $\alpha_{cr}$ discutido na \autoref{sec:ordem-analise} em relação ao
limite normativo.

\section{SITUAÇÕES TRANSITÓRIAS E ACIDENTAIS}\label{sec:verif-transitorias}

Três situações fora do uso corrente foram verificadas.

A primeira é o impacto de veículo em pilar, pertinente por tratar-se de
estacionamento. Sob combinação acidental, com força de 50 kN aplicada a 0,50 m
do piso na direção de marcha, obteve-se $\eta = 0{,}4774$. As seções resistem às
forças de choque de projeto sem alteração. Ainda assim, a proteção física dos
pilares por defensas permanece recomendável, por reduzir dano localizado e custo
de reparo, ainda que não seja exigida pela verificação.

A segunda é o içamento da viga na montagem, verificado conforme a
EN 1991-1-6 \cite{EN1991-1-6:2005}. Com dois pontos de içamento situados a
0,207 do vão, condição que iguala os momentos positivo e negativo, e lingas
inclinadas a 30$^{\circ}$ da vertical, a peça de 75 kN desenvolve momentos de
$-21$ e $+63$ kN$\cdot$m e taxa governante de 0,0420, por cisalhamento. Note-se
que nesta fase a viga não dispõe de contraventamento, e o comprimento efetivo de
flambagem lateral é de 11,9 m, com $k_{crit} = 0{,}797$; ainda assim a taxa
correspondente é de apenas 0,006.

A terceira é a situação de incêndio, verificada pelo método da seção
transversal reduzida do item 4.2.2 da EN 1995-1-2 \cite{EN1995-1-2:2004}, com a
combinação excepcional da EN 1990 \cite{EN1990:2002}. Para o tempo requerido
declarado de 30 minutos, a taxa governante é de 0,2128, no pilar central. A
carbonização reduz a seção resistente, mas a combinação excepcional dispensa os
coeficientes parciais de majoração das ações, e o saldo é favorável nesta
estrutura, cujas peças de madeira operam com folga ampla em situação normal.
Conforme exposto na \autoref{sec:delimitacao}, o tempo requerido é declaração de
projeto, e não prescrição do Eurocódigo.

\section{ESTADOS-LIMITES DE SERVIÇO}\label{sec:verif-els}

A verificação de deformação considera a fluência por meio de $k_{def} = 2{,}00$,
correspondente à classe de serviço 3. A \autoref{tab:verif-els} apresenta os
resultados, obtidos conforme o item 2.3.2.2 e a tabela 7.2 da
EN 1995-1-1 \cite{EN1995-1-1:2004}, com combinações da
EN 1990 \cite{EN1990:2002}. Os limites são parâmetros de determinação nacional,
registrados na \autoref{tab:ndp}.

\tabela{Verificações de estado-limite de serviço}{\footnotesize
\begin{tabular}{lrrcc}
\toprule
    \textbf{Grandeza} & \textbf{Valor} & \textbf{Limite} & \textbf{$\eta$} & \textbf{Situação} \\
\midrule
    Flecha final líquida        & 15,39 mm & 81,0 mm ($L/250$) & 0,1900 & Atende \\
    Flecha instantânea          &  6,56 mm & 67,5 mm ($L/300$) & 0,0972 & Atende \\
    Abertura entre apoios       &  2,56 mm & 40,5 mm ($L/500$) & 0,0632 & Atende \\
\bottomrule
\end{tabular}}{Elaborado pelo autor (2026)}{verif-els}{}{}

Três observações decorrem desses resultados.

A primeira é que a deformação está longe de governar, com taxa máxima de 0,1900.
Isso ocorre apesar de $k_{def} = 2{,}00$ multiplicar por três a parcela
permanente da flecha, e decorre da altura generosa da seção na região central do
vão.

A segunda diz respeito à composição da flecha. A parcela devida ao cisalhamento
responde por 21,2\% do total, valor que a prática corrente frequentemente
despreza. Em peças de madeira, cuja razão entre módulo de elasticidade e módulo
de cisalhamento é da ordem de dezessete, essa parcela não é desprezável, e o
cálculo a incorpora explicitamente.

A terceira é que a combinação característica governante é a de sobrecarga, e não
a de vento, em coerência com o observado para o estado-limite último na
\autoref{sec:esforcos}. A sucção produz, isoladamente, flecha ascendente de
7,39 mm, superior em módulo à flecha descendente de 4,42 mm devida às ações
permanentes. O que decide, porém, é a flecha da combinação: descontada a parcela
permanente, a sucção deixa 2,98 mm de flecha ascendente líquida, contra os
6,56 mm descendentes da combinação com sobrecarga. É esta última que governa.

\section{SÍNTESE}\label{sec:sintese-dimensionamento}

Todas as verificações executadas são atendidas. A taxa de utilização máxima do
conjunto é de 0,8352, correspondente ao contraventamento vertical, e o
veredito global do dimensionamento é de aprovação.

A \autoref{fig:modelo-eta} ilustra a distribuição da taxa de utilização em
flexão ao longo das vigas principais, obtida da tensão normal do modelo de
elementos finitos dividida pela resistência de cálculo à flexão. Ela não
substitui as verificações da \autoref{tab:verif-viga}, que incorporam os
coeficientes de instabilidade e de borda inclinada do Eurocódigo 5 e são as que
governam o dimensionamento.

\figura
{Distribuição da taxa de utilização em flexão nas vigas principais}
{0.98}
{fig/modelo_eta.png}
{Elaborado pelo autor (2026)}
{modelo-eta}
{}
{}

A distribuição das taxas revela um dimensionamento governado pelos sistemas de
estabilização, e não pelos elementos. As cinco verificações mais solicitadas são
os três contraventamentos e as ligações de extremidade das suas diagonais, todas
entre 0,70 e 0,80, seguidas da fixação da telha, com 0,68. A verificação mais
solicitada da viga principal, o cisalhamento, fica em 0,3814, e a menos
solicitada, a flambagem lateral, em 0,1039.

Em uma edificação aberta, de
cobertura leve e grande área exposta, a ação do vento paralelo à cumeeira e a
estabilização dos elementos comprimidos determinam o dimensionamento mais do que
o carregamento gravitacional. A madeira das vigas trabalha com folga; o aço das
diagonais e as fixações da cobertura é que operam próximo dos seus limites.
